%==============================================================================
% Chapitre 40 - Structure d'une cartouche
%==============================================================================

\section{Introduction}

La cartouche moderne regroupe dans un ensemble unique tous les éléments
nécessaires au lancement d'un projectile.

Cette conception constitue l'une des innovations majeures de l'histoire des
armes à feu.

Elle permet :

\begin{itemize}
    \item un chargement rapide ;
    \item une excellente fiabilité ;
    \item une bonne protection des composants ;
    \item une grande reproductibilité des performances.
\end{itemize}

\begin{figure}[htbp]
\centering
\begin{tikzpicture}[scale=1]

% Etui
\draw[thick] (0,0) -- (0,5);
\draw[thick] (2,0) -- (2,5);
\draw[thick] (0,0) -- (2,0);

% Epaulement
\draw[thick] (0,5) -- (0.5,6);
\draw[thick] (2,5) -- (1.5,6);

% Collet
\draw[thick] (0.5,6) -- (0.5,7);
\draw[thick] (1.5,6) -- (1.5,7);

% Projectile
\draw[thick] (0.5,7) -- (1,8);
\draw[thick] (1.5,7) -- (1,8);

% Amorce
\draw[fill=white] (1,0) circle (0.2);

% Labels
\node[right] at (2.3,7.5) {Projectile};
\node[right] at (2.3,3.5) {Poudre};
\node[right] at (2.3,0.2) {Amorce};
\node[left] at (-0.3,3.5) {Étui};

\end{tikzpicture}
\caption{Structure simplifiée d'une cartouche à percussion centrale.}
\label{fig:cartouche_generale}
\end{figure}

\begin{aretenir}

Une cartouche moderne regroupe dans un même ensemble le projectile, l'étui,
la poudre et l'amorce.

\end{aretenir}

\section{Vue d'ensemble}

Une cartouche moderne est généralement constituée de quatre éléments
principaux :

\begin{enumerate}
    \item le projectile ;
    \item l'étui ;
    \item la charge propulsive ;
    \item l'amorce.
\end{enumerate}

Ces éléments fonctionnent ensemble pour produire le tir.

\section{Le projectile}

Le projectile est la seule partie destinée à quitter l'arme.

Son rôle est de transmettre l'énergie vers la cible.

Ses caractéristiques influencent notamment :

\begin{itemize}
    \item la précision ;
    \item la portée ;
    \item la stabilité ;
    \item les effets terminaux.
\end{itemize}

\section{L'étui}

L'étui constitue l'enveloppe de la cartouche.

Il remplit plusieurs fonctions :

\begin{itemize}
    \item maintien des composants ;
    \item protection de la poudre ;
    \item alimentation de l'arme ;
    \item étanchéité lors du tir.
\end{itemize}

\section{La charge propulsive}

La charge propulsive est généralement constituée d'une poudre sans fumée.

Sa combustion génère :

\begin{itemize}
    \item des gaz ;
    \item de la chaleur ;
    \item de la pression.
\end{itemize}

Cette pression accélère le projectile dans le canon.

\section{L'amorce}

L'amorce constitue le dispositif d'allumage.

Lorsqu'elle est frappée par le percuteur :

\begin{itemize}
    \item elle produit une flamme ;
    \item cette flamme enflamme la poudre ;
    \item la combustion débute.
\end{itemize}

\section{Fonctionnement général}

La séquence simplifiée d'un tir est la suivante :

\begin{enumerate}
    \item percussion de l'amorce ;
    \item inflammation de la poudre ;
    \item production de gaz ;
    \item augmentation de la pression ;
    \item mouvement du projectile ;
    \item sortie du projectile du canon.
\end{enumerate}

Cette séquence sera étudiée en détail dans la partie consacrée à la
balistique intérieure.

\section{Cartouche à percussion centrale}

La majorité des cartouches modernes utilisent une amorce située au centre du
culot.

On parle alors de percussion centrale.

Cette architecture présente :

\begin{itemize}
    \item une grande robustesse ;
    \item une bonne fiabilité ;
    \item une forte résistance à la pression.
\end{itemize}

\section{Cartouche à percussion annulaire}

Certaines cartouches utilisent une amorce répartie dans le bourrelet de
l'étui.

On parle alors de percussion annulaire.

Le cas le plus connu est celui du :

\begin{itemize}
    \item .22 Long Rifle.
\end{itemize}

Ces cartouches sont généralement destinées à des pressions plus faibles.

\section{Importance de l'assemblage}

Les performances d'une cartouche dépendent non seulement de ses composants,
mais également de leur assemblage.

Les paramètres importants incluent notamment :

\begin{itemize}
    \item la longueur totale ;
    \item le sertissage ;
    \item la masse de poudre ;
    \item l'enfoncement du projectile.
\end{itemize}

\section{Vers l'étude détaillée des composants}

La compréhension complète d'une cartouche nécessite maintenant l'étude de
chacun de ses éléments.

Nous commencerons par l'étui, composant souvent discret mais pourtant
essentiel au fonctionnement et à la sécurité du système.

\section{À retenir}

\begin{aretenir}

\begin{itemize}
    \item Une cartouche moderne comprend quatre éléments principaux.
    \item L'amorce initie le tir.
    \item La poudre fournit l'énergie.
    \item L'étui assure le maintien et l'étanchéité.
    \item Le projectile est la partie destinée à quitter l'arme.
\end{itemize}

\end{aretenir}
